\documentclass[11pt,a4paper]{article}
\usepackage{amsmath,amssymb,amsthm}
\usepackage[margin=2.4cm]{geometry}
\newtheorem*{theorem}{Theorem}
\newtheorem{lemma}{Lemma}
\newcommand{\abs}[1]{\lvert #1\rvert}
\emergencystretch=1.5em
\title{\large\bfseries The bound $A(18,6,6)\le 183$}
\author{Jiang Chang}
\date{November 2026}
\begin{document}
\maketitle

\paragraph{Notation.}
$A(n,d,w)$ is the maximum size of a binary code of length $n$, constant
weight $w$ and minimum distance at least $d$. Throughout, $n=18$, $d=6$ and
$w=6$. A \emph{word} is a $6$-subset of $[18]=\{1,\dots,18\}$; there are
$V=\binom{18}6=18564$ words. Two words $u,u'$ are at distance
$2(6-\abs{u\cap u'})$. So a \emph{code} is a set of words any two of which
meet in at most $3$ points. A set of words is called \emph{compatible} if it
is a code.

\begin{theorem}
$A(18,6,6)\le 183$.
\end{theorem}

The bounds known before were $133\le A(18,6,6)\le 186$, see~\cite{Btab}.
The upper bound $186$ follows from $A(17,6,6)\le 124$ (\"Osterg{\aa}rd
\cite{Ost}) and the Johnson bound $A(18,6,6)\le\lfloor 18\cdot 124/12\rfloor$.

\paragraph{Outline.}
The proof is a semidefinite programming bound. Sections 1--6 give a list of
linear and semidefinite constraints on certain numbers $z(S)$ attached to a
code, and prove that every code satisfies them. Section~7 describes a
certificate, in rational numbers, that these constraints imply $\abs C<184$;
it is checked by computer in exact arithmetic. The constraints are those of
the four-point bound of Polak~\cite{Polak} together with further
constraints obtained by counting. Adding such constraints to a bound of
this kind goes back to~\cite{BB} and~\cite[Theorem 22]{BBMOS}. Among them
are $23$ inequalities that come from the leaves of packings of $4$-sets
(Sections 3--5); such leaves are classical, see~\cite{B79}.

\paragraph{What is proved, and how it is checked.}
The Theorem rests on three things: (i) the constraints are valid for every
code, which is proved in Sections 1--6; (ii) the certificate of Section~7,
which is verified by the public program in exact arithmetic; (iii) the
$16387$ matrices in the public model file are correct instances of these
constraints. For (iii), every matrix in the model carries a label naming
its family and its parameters, and a written specification states how the
coefficients follow from the label. All matrices have been computed again
from these definitions, in exact integer arithmetic, by programs written
independently of the program that generated the model, and compared with
the model coefficient by coefficient: the $23$ inequalities of Section~5 by
the public program; the $269$ matrices of Section~2, in the form reduced
by symmetry, by one reconstruction; and the remaining $16095$ constraints,
namely the $941$ inequalities $z(S)\ge 0$ and the $15154$ constraints of
Section~6, by another. Together these
account for all $2074973$ coefficients in the model. See Section~7 for the
files.

\section{The numbers $z(S)$ and two principles}

For a word $u$ and a code $C$ put $X_u(C)=1$ if $u\in C$, and $X_u(C)=0$
otherwise. For a finite set $S$ of words put $X_S=\prod_{u\in S}X_u$, with
$X_\emptyset=1$. Then $X_SX_{S'}=X_{S\cup S'}$, and $X_S$ vanishes on every
code if $S$ is not compatible. A \emph{polynomial} is a real linear
combination $p=\sum_S p_SX_S$; its \emph{degree} is the largest $\abs S$
with $p_S\ne 0$. It is a function on the set of all codes.

Let $C$ be a code. For a set $S$ of at most four words put
\[
z(S)=\frac{V}{18!}\sum_{g}X_S(gC),
\]
where $g$ runs over the permutations of $[18]$. Then $z(S)$ depends only on
the orbit of $S$ under these permutations, and
\[
z(\emptyset)=V,\qquad z(\{u\})=\abs C,\qquad 0\le z(S)\le V,\qquad
z(S)=0\ \text{if $S$ is not compatible.}
\]
(For the second: the average of $X_u(gC)$ over $g$ is $\abs C/V$.)

\begin{itemize}
\item[(A)] Let $p$ be a polynomial of degree at most $4$ such that
$p(C')\ge 0$ for every code $C'$. Then $\sum_Sp_S\,z(S)\ge 0$.
\item[(B)] Let $h$ be a polynomial with $h(C')\ge 0$ for every code $C'$,
and let $f_1,\dots,f_r$ be polynomials such that every product $hf_if_j$
has degree at most $4$. Then the $r\times r$ matrix with entries
$\sum_S(hf_if_j)_S\,z(S)$ is positive semidefinite.
\end{itemize}
Indeed, $\sum_Sp_Sz(S)=\frac{V}{18!}\sum_gp(gC)$, and every $gC$ is a code.
For (B), the quadratic form of the matrix at a vector $a$ is
$\frac V{18!}\sum_g h(gC)\bigl(\sum_ia_if_i(gC)\bigr)^2\ge 0$. Statement (B)
is the principle of moment and localizing matrices~\cite{Las}.

The unknowns of the semidefinite program are the numbers $z(\omega)$, one
for each orbit $\omega$ of compatible sets of one to four words; there are
$941$ of them. Every constraint below is an instance of (A) or (B). A
principal submatrix of a positive semidefinite matrix, and its restriction
to an invariant subspace of a group of symmetries, are again positive
semidefinite; this is used to reduce the sizes of the matrices.

\section{The constraints of the four-point bound}

Let $D$ be a compatible set of words, and consider the sets $S\supseteq D$
with $\abs D+2\abs{S\setminus D}\le 4$. For two such sets $S,S'$ we have
$X_DX_SX_{S'}=X_{S\cup S'}$, of degree at most $4$. By (B), with $h=X_D$ and
the polynomials $X_S$, the matrix with entries $z(S\cup S')$ is positive
semidefinite. These matrices, and the inequalities $z(S)\ge 0$, are the
constraints of~\cite{Polak}.

\section{Local packings and their leaves}

For disjoint sets $R,Z\subseteq[18]$ the \emph{interval} $[R,Z]$ is the set
of words $u$ with $R\subseteq u$ and $u\cap Z=\emptyset$. Deleting the
points of $R$ and $Z$ turns the words of a code in $[R,Z]$ into subsets of
size $6-\abs R$ of a set of $18-\abs R-\abs Z$ points, any two at distance
at least $6$.

In this section and the next two, $\abs R=2$ and $\abs Z=b$. Let
$Y=[18]\setminus(R\cup Z)$ and $m=\abs Y=16-b$. For a code $C$, the sets
$u\setminus R$ with $u\in C\cap[R,Z]$ are $4$-subsets of $Y$, called
\emph{blocks}. Two blocks have at most one common point, because the two
words have at most three. So the blocks form a \emph{packing} $P$: every
pair of points of $Y$ lies in at most one block. For a $4$-subset $B$ of
$Y$ we write $X_B$ for $X_{R\cup B}$; a set of $4$-subsets of $Y$ is
compatible if it is a packing.

Assume that $m\ge 4$ and $m\equiv 1$ or $4\pmod{12}$, that is, $b\in\{0,3,12\}$ and
$m\in\{16,13,4\}$. Put
\[
K=\frac{m(m-1)}{12},\qquad N=\abs P,\qquad t=K-N .
\]
The \emph{leave} $L$ of $P$ is the graph on $Y$ whose edges are the pairs
that lie in no block. A point in $r$ blocks has degree $m-1-3r$ in $L$, and
$L$ has $\binom m2-6N=6t$ edges. So $t$ is a nonnegative integer, and
every degree of $L$ is divisible by $3$.

\begin{lemma}
If $t=0$ then $L$ has no edge. If $t=1$ then $L$ is a complete graph on
four points, together with isolated points.
\end{lemma}

\begin{proof}
Let $t=1$. Then $L$ has $6$ edges and degree sum $12$. Every non-isolated
point has degree at least $3$, so there are at most four of them; and $6$
edges need at least four points.
\end{proof}

Let $T$ be the number of triangles of $L$. Let $I$ be the number of
unordered pairs of blocks with exactly one common point, and $Q$ the number
of unordered triples of blocks that pairwise have exactly one common point
and have no point common to all three.

\begin{lemma}
$T=\binom m3-(6m-20)\,N+9I-Q$.
\end{lemma}

\begin{proof}
A triple of points is a triangle of $L$ if none of its three pairs lies in a
block. By inclusion and exclusion, $T=\binom m3-\sigma_1+\sigma_2-\sigma_3$,
where $\sigma_i$ counts the choices of a triple and of $i$ of its pairs that
all lie in blocks. There are $6N$ pairs in blocks, each in $m-2$ triples, so
$\sigma_1=6(m-2)N$. Two pairs of a triple lie either in the same block,
which then contains the triple ($3\cdot 4N$ choices), or in two blocks with
exactly one common point, the common point of the two pairs ($3\cdot3$
choices for each such pair of blocks); so $\sigma_2=12N+9I$. If all three
pairs lie in blocks, the blocks are either the same ($4N$ triples) or three
different ones, as counted by $Q$; two of the pairs cannot lie in one block
and the third in another, since the first block contains the triple and the
third pair would lie in two blocks. Three blocks as in the definition of
$Q$ meet pairwise in three different points, which form one such triple.
So $\sigma_3=4N+Q$.
\end{proof}

\begin{lemma}
$F:=T+4t(t-2)\ge 0$ for every packing $P$.
\end{lemma}

\begin{proof}
If $t=0$ then $T=0$ and $F=0$. If $t=1$ then $T=4$ by Lemma~1, and $F=0$.
If $t\ge 2$ both terms are nonnegative.
\end{proof}

By Lemma~2, and since $N=\sum_BX_B$ over all $4$-subsets $B$ of $Y$ and
$N^2=N+2\sum_{\{B,B'\}}X_BX_{B'}$, the function $F$ is a polynomial of
degree at most $3$ in the indicators of the blocks, $F=\sum_Sf(S)X_S$ over
the compatible sets $S$ of blocks, with
\begin{center}
\begin{tabular}{ll}
\hline
Set $S$ of blocks & $f(S)$ \\
\hline
empty & $\binom m3+4K(K-2)$ \\
one block & $32-8K-6m$ \\
two blocks & $17$ if they have one common point, $8$ if they are disjoint\\
three blocks & $-1$ if they are as in the definition of $Q$, otherwise $0$\\
four or more blocks & $0$\\
\hline
\end{tabular}
\end{center}
The following products are also nonnegative on every code, and have degree
at most $4$:
\[
N\!F,\qquad tF=(K-N)F,\qquad H_{j,h}F ,
\]
where, for $0\le j\le 2$ and $0\le h\le b$, $H_{j,h}$ is the number of
words $u$ of the code with $\abs{u\cap R}=j$ and $\abs{u\cap Z}=h$. Since
$H_{2,0}=N$, the pair $(j,h)=(2,0)$ is omitted. The coefficients of $N\!F$
are $g(S)=\abs S\,f(S)+\sum_{B\in S}f(S\setminus\{B\})$, because
$X_BX_S$ is $X_S$ if $B\in S$ and $X_{S\cup\{B\}}$ otherwise. The
coefficients of $tF$ are $K\,f(S)-g(S)$.

\section{A lemma on graphs with twelve edges}

A \emph{$4$-cycle} of a graph is a subgraph that is a cycle of length $4$,
induced or not.

\begin{lemma}
Let $\Gamma$ be a graph with $12$ edges in which every degree is divisible
by $3$. Let $T$ be its number of triangles and $C_4$ its number of
$4$-cycles. Then $T+C_4\ge 4$.
\end{lemma}

\begin{proof}
Delete the isolated points; let $s$ points remain. Their degrees are at
least $3$ and add up to $24$, so $s\le 8$; hence the degrees are at most
$7$, that is, they are $3$ or $6$. If $s\le 6$, all degrees are at most
$5$, hence equal to $3$, and their sum is at most $18$. So $s=7$ or $s=8$.

\emph{The case $s=7$.} One point $x$ has degree $6$ and the others have
degree $3$. Each of the other six points has two neighbours other than $x$,
so there are six edges not on $x$, and each forms a triangle with $x$.
Hence $T\ge 6$.

\emph{The case $s=8$.} Every degree is $3$. For two points $x,y$ let
$c(x,y)$ be the number of their common neighbours. Counting paths of length
two; triangles, by their edges; and $4$-cycles, by their two pairs of
opposite points, gives
\[
\sum_{\{x,y\}}c(x,y)=8\binom32=24,\qquad
\sum_{\{x,y\}\ \text{an edge}}c(x,y)=3T,\qquad
\sum_{\{x,y\}}\binom{c(x,y)}2=2C_4 .
\]
There are $\binom82-12=16$ pairs that are not edges, and
$\binom c2\ge c-1$ for every integer $c\ge 0$. Adding this over the
non-edges gives
\[
2C_4\ \ge\ (24-3T)-16=8-3T .
\]
So $C_4\ge 4$ if $T=0$, and $C_4\ge 3$ if $T=1$. If $T\ge 4$ there is
nothing to prove. It remains to consider $T=2$ and $T=3$.

Two triangles with a common point have a common edge, since the common
point has only three neighbours. Let $xyp$ and $xyq$ be triangles. Then $x$
and $y$ have no further neighbours. If $p$ and $q$ are adjacent, the four
points carry four triangles. Otherwise $p$ and $q$ each have one neighbour
among the other four points. These four points have degree sum $12$, of
which $2$ is used by the edges to $p$ and $q$; so they carry $5$ edges,
which form two triangles with a common edge. Again $T\ge 4$.

So if $T\le 3$ the triangles are pairwise disjoint, and there are at most
two of them, as there are only $8$ points. Let $T=2$, with triangles $A$ and
$B$ and remaining points $x,y$. A point outside a triangle is adjacent to
at most one of its points, since otherwise there is a third triangle. So
$x$ has at most one neighbour in $A$ and at most one in $B$; having three
neighbours, it is adjacent to $y$ and to exactly one point of $A$, say
$a$, and one of $B$. In the same way $y$ is adjacent to exactly one point
$a'$ of $A$. Here $a'\ne a$, since $a$ has only one neighbour outside $A$.
Now $x,a,a',y$ form a $4$-cycle. There is another one through $B$. So
$C_4\ge 2$ and $T+C_4\ge 4$.
\end{proof}

For example, the graph on the points $0,\dots,7$ with edges $\{i,i+1\}$
and $\{i,i+4\}$, taken modulo $8$, has $T=0$ and $C_4=4$. The case $s=8$ was also confirmed by computer, by listing all
$19355$ graphs on $8$ given points in which every degree is $3$.

\begin{lemma}
Let $P$ be a packing as in Section~$3$, and let $\Sigma=T+C_4$ for its
leave. Then
\[
H_2:=\Sigma+5t^2-12t\ \ge 0,\qquad
H_4:=24\,\Sigma+t(t-3)(t-5)(5t-26)\ \ge 0 .
\]
\end{lemma}

\begin{proof}
If $t=0$ the leave is empty and $\Sigma=0$. If $t=1$ the leave is a complete
graph on four points, with $4$ triangles and $3$ $4$-cycles, so $\Sigma=7$.
If $t=2$ the leave has $12$ edges and all degrees divisible by $3$, so
$\Sigma\ge 4$ by Lemma~4. Always $\Sigma\ge 0$. The values of the two
polynomials in $t$ are
\begin{center}
\begin{tabular}{lcccccc}
\hline
$t$ & $0$ & $1$ & $2$ & $3$ & $4$ & $5$ \\
\hline
$5t^2-12t$ & $0$ & $-7$ & $-4$ & $9$ & $32$ & $65$ \\
$t(t-3)(t-5)(5t-26)$ & $0$ & $-168$ & $-96$ & $0$ & $24$ & $0$ \\
\hline
\end{tabular}
\end{center}
and both are positive for $t\ge 6$.
\end{proof}

For a pair $e$ of points of $Y$ let $Y_e=1-\sum_{B\ni e}X_B$. On a packing
it is $1$ if $e$ is an edge of the leave and $0$ otherwise. So $T$ is the
sum, over the triangles of the complete graph on $Y$, of the products of
three $Y_e$, and $C_4$ is the corresponding sum over its $4$-cycles, with
four factors. Hence $H_2$ and $H_4$ are polynomials of degree at most $4$
in the indicators of the blocks. Their coefficients are obtained from their
values: for a compatible set $S$ of at most four blocks,
\[
c(S)=\sum_{S'\subseteq S}(-1)^{\abs{S\setminus S'}}\,H(S'),
\]
where $H(S')$ is the value of $H_2$ or $H_4$ on the packing $S'$ in the
same set $Y$. (As $H(P)=\sum_{S\subseteq P}c(S)$ for every packing $P$,
this is M\"obius inversion on the subsets of $S$; all subsets of a
compatible set are compatible.)

\section{From one interval to the numbers $z(S)$}

Fix $b$, and let $p_{R,Z}$ be one of the polynomials $F$, $N\!F$, $tF$,
$H_{j,h}F$, $H_2$, $H_4$ for the interval $[R,Z]$, written in the
indicators $X_u$ of the words. It is nonnegative on every code. So is the
sum $p=\sum p_{R,Z}$ over all pairs $(R,Z)$, and by (A)
\begin{equation}
\sum_Sp_S\,z(S)\ \ge 0. \tag{1}
\end{equation}
The coefficients $p_S$ are found as follows. For a nonempty set $S$ of
words let $\iota(S)$ be the number of points in all words of $S$, and
$\upsilon(S)$ the number of points in at least one. Then $S$ lies in
$\binom{\iota(S)}2\binom{18-\upsilon(S)}b$ intervals (here and below
$\binom xy=0$ if $y<0$ or $y>x$): $R$ consists of two points lying in all
words of $S$, and $Z$ of $b$ points lying in none. Any two such choices
differ by a permutation of the points that fixes every word of $S$, so the
corresponding sets of blocks have the same coefficients. So for $F$,
$N\!F$, $tF$, $H_2$ and $H_4$,
\[
p_S=c(S)\binom{\iota(S)}2\binom{18-\upsilon(S)}b,\qquad
p_\emptyset=c(\emptyset)\binom{18}2\binom{16}b ,
\]
where $c$ denotes the coefficients found in Sections 3 and 4.

For $H_{j,h}F$, a term arises from a word $u$ counted by $H_{j,h}$ and a
set $S'$ of at most three words in the interval, with coefficient
$f(S')$; here $u\notin S'$ because $(j,h)\ne(2,0)$. If $S'$ is not empty,
let $A$ be the set of points in all words of $S'$ and $B$ the set of points
in none. The number of intervals containing $S'$ and with
$\abs{u\cap R}=j$, $\abs{u\cap Z}=h$ is
\[
\binom{\abs{A\cap u}}j\binom{\abs{A\setminus u}}{2-j}
\binom{\abs{B\cap u}}h\binom{\abs{B\setminus u}}{b-h},
\]
and if $S'$ is empty it is
$\binom6j\binom{12}{2-j}\binom{6-j}h\binom{10+j}{b-h}$. The coefficient
of $X_S$ is the sum, over the words $u\in S$, of $f(S\setminus\{u\})$
times this number.

Finally (1) is written in the unknowns $z(\omega)$. All sets $S$ in an
orbit $\omega$ have the same coefficient, so $z(\omega)$ gets the
coefficient $p_S$ times the size of $\omega$. For a set of $r$ distinct words, numbered in some order, let
$h_\alpha$ be the number of points whose incidence with the $r$ words is
the pattern $\alpha\in\{0,1\}^r$, and let $\sigma$ be the number of
permutations of the $r$ words that preserve all the numbers $h_\alpha$.
The permutations of the points that map the set to itself number
$\sigma\prod_\alpha h_\alpha!$, so the orbit has
$18!/(\sigma\prod_\alpha h_\alpha!)$ elements. The constant term is
$p_\emptyset V$, as $z(\emptyset)=V$. Each inequality is divided by the
greatest common divisor of its coefficients.

\paragraph{The $23$ inequalities.}
\begin{center}
\begin{tabular}{rrrp{11.2cm}}
\hline
$b$ & $m$ & $K$ & Inequalities \\
\hline
$0$ & $16$ & $20$ & $F$, $N\!F$, $tF$; $H_{j,h}F$ for $(j,h)=(0,0),(1,0)$;
 $H_2$, $H_4$: seven inequalities\\
$3$ & $13$ & $13$ & $F$, $N\!F$, $tF$; $H_{j,h}F$ for the eleven pairs
 $(j,h)\ne(2,0)$ with $0\le j\le 2$, $0\le h\le 3$; $H_2$, $H_4$: sixteen
 inequalities\\
$12$ & $4$ & $1$ & none: all these polynomials vanish identically\\
\hline
\end{tabular}
\end{center}
For $m=4$ there is at most one block, the leave is complete or empty, and
$F$, $H_2$, $H_4$ are zero. The same construction for the complements of
the words (weight $12$) requires $\abs R=8$ to leave blocks of size $4$, so
that $m\le 10$, and then $m=4$.

\section{The other constraints}

In this section $[R,Z]$ is an arbitrary interval, $k=6-\abs R$,
$m=18-\abs R-\abs Z$, and $N$ is the number of words of the code in
$[R,Z]$. Then $\binom Nr=\sum X_S$ over the $r$-sets $S$ of words in
$[R,Z]$, so a polynomial of degree $r$ in $N$ is a polynomial of degree at most $r$
in the indicators. The components of the vectors $f$ below are polynomials of
degree at most $1$: the constant $1$, counts of the words of the code in
given classes with respect to a fixed partition of the points, and integer
linear combinations of the indicators $X_u$ that arise when the matrices
are reduced by symmetry. In each family only products of degree at most
$4$ are used, and the polynomials or matrices for all intervals with the
same $\abs R$ and $\abs Z$ are added, as in Section~5. An interval with
$R=\emptyset$ consists of the words disjoint from $Z$, that is, of the
words whose complements contain $Z$.

\emph{Sets that meet a code in at most one word.} Let $r\ge 0$, let $E$ be a set of
$4+2r$ points and let $\mathcal Q$ be the set of words that have at least
$4+r$ points in $E$. Two words of $\mathcal Q$ have at least
$2(4+r)-(4+2r)=4$ common points. So a code contains at most one word of
$\mathcal Q$, and $q=1-\sum_{u\in\mathcal Q}X_u\ge 0$ on every code. The
constraints are (A) for $X_Dq$ with $\abs D\le 3$, and (B) for $h=X_uq$.

\emph{Intervals with pairwise disjoint blocks.} If $\min(k,m-k)=3$, the
words in $[R,Z]$, or their complements in $Y$, are pairwise disjoint after
deleting $R$. If moreover $9\le m<12$, then $N\le 3$, and for
$j=0,1,2,3$
\[
\sum_{r=j}^3(-1)^{r-j}\binom rj\binom Nr
\]
is $1$ if $N=j$ and $0$ otherwise. The constraints are (A) for the product
of $X_u$ and this polynomial.

\emph{Polynomials in $N$.} Let $J(m,k)=1$ if $\min(k,m-k)<3$, and otherwise
let $J(m,k)$ be the smaller of $\lfloor m\,J(m-1,k-1)/k\rfloor$ and
$\lfloor m\,J(m-1,k)/(m-k)\rfloor$. The words on a given point, and the words not on it, form codes to which
induction applies. Counting the pairs (word, point on it) and (word, point
not on it) then gives $kN\le m\,J(m-1,k-1)$ and $(m-k)N\le m\,J(m-1,k)$,
so $N\le J(m,k)$. If
$\pi$ is a real polynomial of degree at most $4$ with $\pi(i)\ge 0$ for
all integers $0\le i\le J(m,k)$, then $\pi(N)\ge 0$ on every code. The
constraints are (A) for such $\pi(N)$.

\emph{Matrices from $N$.} For an integer $j$, $(N-j)(N-j-1)\ge 0$, and
$N(J(m,k)-N)\ge 0$. The constraints are (B) with $h$ one of these.

\emph{The known bound.} Let $M=\sum_uX_u$ be the number of words of the
code. Then $186-M\ge 0$ on every code. The constraints are (A) for
$X_D(186-M)$ with $\abs D\le 3$, and (B) for $h=X_u(186-M)$. The number
$183$ is not used anywhere.

\section{The certificate}

Number the unknowns $z_1,\dots,z_{941}$, with $z_1=z(\{u\})=\abs C$. Each
constraint has the form: the symmetric matrix
$F_\beta(z)=F_{\beta,0}+\sum_iz_iF_{\beta,i}$ is positive semidefinite,
where the $F_{\beta,i}$ are integer matrices; a linear inequality is the
case of a $1\times 1$ matrix. There are $16387$ such matrices, the largest
of order $47$. The certificate consists of rational positive semidefinite
matrices $Y_\beta$ of the same orders. With
$\langle A,B\rangle=\sum_{i,j}A_{ij}B_{ij}$, put
\[
c_0=\sum_\beta\langle F_{\beta,0},Y_\beta\rangle,\qquad
\rho_i=\sum_\beta\langle F_{\beta,i},Y_\beta\rangle+
\begin{cases}1&(i=1)\\0&(i\ne1).\end{cases}
\]
For every code, $\langle F_\beta(z),Y_\beta\rangle\ge 0$ for all $\beta$,
as both matrices are positive semidefinite.
Adding, $c_0-z_1+\sum_i\rho_iz_i\ge 0$. Since $0\le z_i\le V=18564$,
\[
\abs C=z_1\ \le\ c_0+18564\sum_i\max(\rho_i,0)=:U .
\]
The exact values are
\begin{align*}
c_0&=183.97846380536651084584137185081,\\
18564\textstyle\sum_i\max(\rho_i,0)&=\phantom{00}0.0003964175429246453124659046,\\
U&=183.97886022290943549115383775541 .
\end{align*}
These are rational numbers with denominators dividing $10^{29}$, written
out in full. As $\abs C$ is an integer and $U<184$, the Theorem follows.
Bounding the terms $\rho_iz_i$ in this way is a standard device for exact
certificates, see the appendix of~\cite{BP}.

\paragraph{Files and programs.}
All files are in the following repository, or attached to its releases.
\begin{center}
\begin{tabular}{ll}
Repository: & \texttt{github.com/wille614/cw-code-upper-bounds}\\
The matrices $F_{\beta,i}$, with their labels: & \texttt{A18\_6\_6\_model.json.gz}\\
The matrices $Y_\beta$: & \texttt{A18\_6\_6\_certificate.json.gz}
\end{tabular}
\end{center}
In a constraint, a constant $a$ in
a polynomial appears in $F_{\beta,0}$ as $aV$, in accordance with
$z(\emptyset)=V$.

The command \texttt{python verify.py 18 6 6} computes the coefficients of
the $23$ inequalities of Section~5 again, from direct counts of triangles
and $4$-cycles and of intervals, and compares them with the model; it
verifies that every $Y_\beta$ is positive semidefinite, by exact
elimination in rational arithmetic; and it computes $c_0$, all $\rho_i$ and
$U$ exactly.

The directory \texttt{reconstruction/A18\_6\_6/} contains the specification
of the labels and two sets of programs with their reports. One computes
again the $941$ inequalities $z(S)\ge 0$ and the $15154$ constraints of
Section~6, the other the $269$ matrices of Section~2 in the form reduced by
symmetry.
Each compares its result with the model, coefficient by coefficient.

\begin{thebibliography}{9}\small
\bibitem{BB} M. R. Best and A. E. Brouwer, The triply shortened binary
Hamming code is optimal, \emph{Discrete Math.} \textbf{17} (1977) 235--245.
\bibitem{BBMOS} M. R. Best, A. E. Brouwer, F. J. MacWilliams, A. M. Odlyzko
and N. J. A. Sloane, Bounds for binary codes of length less than 25,
\emph{IEEE Trans. Inform. Theory} \textbf{24} (1978) 81--93.
\bibitem{B79} A. E. Brouwer, Optimal packings of $K_4$'s into a $K_n$,
\emph{J. Combin. Theory Ser. A} \textbf{26} (1979) 278--297.
\bibitem{Btab} A. E. Brouwer, Bounds for binary constant weight codes,
\texttt{https://aeb.win.tue.nl/codes/Andw.html}.
\bibitem{BP} A. E. Brouwer and S. C. Polak, Uniqueness of codes using
semidefinite programming, \emph{Des. Codes Cryptogr.} \textbf{87} (2019)
1881--1895.
\bibitem{Las} J. B. Lasserre, Global optimization with polynomials and the
problem of moments, \emph{SIAM J. Optim.} \textbf{11} (2001) 796--817.
\bibitem{Ost} P. R. J. \"Osterg{\aa}rd, Classification of binary constant
weight codes, \emph{IEEE Trans. Inform. Theory} \textbf{56} (2010)
3779--3785.
\bibitem{Polak} S. C. Polak, Semidefinite programming bounds for
constant-weight codes, \emph{IEEE Trans. Inform. Theory} \textbf{65}
(2019) 28--38.
\end{thebibliography}
\end{document}
